\chapter*{Introducción} \label{resumen}
\addcontentsline{toc}{chapter}{Introducción}

El cómputo con unidades de procesamiento gráfico (GPUs) se ha consolidado como una de las vías principales para acelerar cargas de trabajo de alto coste computacional, tal y como puede verse en ámbitos tan destacados como el aprendizaje automático o la simulación científica \cite{ibm_hpc}. El modelo habitual de despliegue consiste en adquirir una o varias tarjetas de una misma familia y generación, gestionadas por un \textit{vendor} de software (driver, \textit{runtime}) homogéneo. No obstante,
la constante demanda de mejora de rendimiento en los sistemas fomenta recurrir a la compra de hardware cada vez
 más nuevo, dejándose atrás un parque de dispositivos muy amplio compuesto por GPUs de años, generaciones y fabricantes 
 muy distintos. Ello es consecuencia directa de la rápida obsolescencia de las tarjetas, unida a la tendencia de 
 reciclar equipos entre otros factores.

Surge entonces una pregunta natural: \emph{¿es posible aprovechar todas esas tarjetas a la vez dentro de una misma máquina?} La respuesta intuitiva es que sí, pero la práctica revela que existen numerosas restricciones técnicas ---a nivel de driver, de arquitectura y de virtualización--- que condicionan e incluso impiden dicha integración. Este trabajo se plantea como un estudio experimental para responder a esa pregunta de forma rigurosa: construiremos un sistema real con varias GPUs de características diferentes y documentaremos, paso a paso, una serie de pruebas explorando qué es posible hacer con ellas y qué no.

El caso de estudio dispone de tres tarjetas NVIDIA de la arquitectura GT200 (2008--2009): dos \textbf{NVIDIA Tesla C1060} (tarjetas de cómputo, \textit{compute capability} 1.3) y una \textbf{ASUS ENGTS250} (tarjeta gráfica de consumo, \textit{compute capability} 1.1). La configuración experimental empleó dos GPUs, una Tesla C1060 y la ENGTS250, porque la tercera tarjeta no fue detectada a través del \textit{riser} PCIe. Ambos modelos están limitados a la última rama de driver NVIDIA denominada \textit{legacy}, 
versión 340.108, cuyo último \textit{release} oficial data de 2019. Este detalle, aparentemente anecdótico, condiciona todo el trabajo y lo convierte en un excelente banco de pruebas para estudiar la viabilidad de sistemas multi-GPU heterogéneos con limitaciones tecnológicas especialmente remarcables.

\section*{Objetivos}
\addcontentsline{toc}{section}{Objetivos}

El propósito general de este trabajo es doble: estudiar experimentalmente la viabilidad de 
aprovechar múltiples GPUs heterogéneas en un mismo sistema mediante un caso de estudio real, 
y sistematizar los criterios de viabilidad empleados en dicho estudio en un marco teórico 
aplicable \textit{a priori} a otras combinaciones de hardware. La manera de proceder para la construcción del marco teórico 
ha sido de forma 
inductiva: se formula y se completa a partir de los propios obstáculos que la experimentación 
ha ido revelando, en lugar de imponerse como un criterio fijado de antemano. Este segundo objetivo 
se desglosa en los siguientes:
\begin{itemize}
    \item Diseñar y construir un sistema que integre varias GPUs de generaciones y propósitos distintos,
     identificando las posibles configuraciones que lo hacen posible.

    \item Analizar la limitación que supone disponer de un único driver para un sistema multi-GPU y evaluar los dos escenarios que 
    derivan de esta condición: emplear un mismo driver para todas las tarjetas (opción no siempre disponible), o aislar cada tarjeta en un entorno 
    propio con su versión de driver y su sistema operativo correspondiente.

    \item Instalar y validar un driver \textit{legacy} de NVIDIA sobre un sistema con 
    kernel y \textit{toolchain} modernos, resolviendo las incompatibilidades que surjan 
    entre ambos.

    \item Proporcionar acceso a las GPUs desde un entorno aislado mediante el paso de 
    dispositivos, e instalar el toolkit CUDA compatible con las arquitecturas disponibles.

    \item Someter a pruebas las tarjetas para medición y comparación de rendimientos (cómputo y ancho de banda de memoria), así como evaluar la 
    contención al ejecutar cargas de trabajo concurrentes.

    \item Documentar los fallos y limitaciones encontrados, tanto de hardware como de 
    software, y plasmarlos en una serie de parámetros de estudio en un marco teórico de criterios de viabilidad  
    generalizable para otros sistemas de características similares.
\end{itemize}

\section*{Motivación}
\addcontentsline{toc}{section}{Motivación}

La motivación de este trabajo es doble. En primer lugar, existe una motivación económica y de sostenibilidad: 
la adquisición de hardware de cómputo nuevo es costosa y el ciclo de vida útil de las tarjetas gráficas es corto,
por lo que reutilizar tarjetas de segunda mano o rescatadas de otros equipos resulta atractivo tanto por su bajo
coste como por reducir el residuo electrónico. Si se demostrase que es posible combinar tarjetas de distintas 
generaciones y marcas en una misma máquina, cualquier usuario podría construir un nodo de cómputo aprovechando 
hardware que de otro modo se descartaría.

En segundo lugar, el problema tiene un interés técnico genuino. Las GPUs no son componentes fácilmente intercambiables: 
dependen de un driver propietario, de arquitecturas de cómputo concretas y de un \textit{runtime} 
(CUDA en el caso de nuestras GPUs NVIDIA). La heterogeneidad introduce
además complejidad en el aislamiento, la gestión de memoria y la contención de recursos. Estudiar dónde están 
los límites de esta heterogeneidad, y cómo pueden salvarse con las distintas herramientas actuales, 
es un problema abierto que combina conocimientos de sistemas operativos, arquitectura 
de computadores y virtualización.

Como se comprobará a lo largo de la memoria, el trabajo también se justifica por la gran cantidad de problemas
reales ---y a menudo no documentados--- que surgen al combinar componentes \textit{legacy} con software y hardware
modernos. Cada uno de esos problemas constituye un resultado en sí mismo.

\section*{Metodología de trabajo}
\addcontentsline{toc}{section}{Metodología de trabajo}

Dado el carácter experimental del trabajo, la metodología seguida es la propia de una bitácora de experimentación estructurada. En cada fase se ha procedido de la siguiente manera:

\begin{enumerate}
    \item \textbf{Planteamiento}: se describe el escenario que se quiere conseguir y se estudian las 
    alternativas técnicas disponibles.

    \item \textbf{Experimentación}: se implementa el escenario sobre el hardware real y se documentan los obstáculos encontrados, con su causa raíz cuando es posible determinarla.

    \item \textbf{Validación}: cada fase concluye con una comprobación objetiva del resultado.

    \item \textbf{Discusión}: se analizan los resultados obtenidos y su relevancia para el objetivo general de viabilidad.
\end{enumerate}

En los casos en los que una vía resulta en descarte, 
el fracaso se documenta y se utiliza como justificación técnica de la opción alternativa, siguiendo la dinámica de que en un 
trabajo experimental los resultados negativos son tan informativos como los positivos.


